\begin{abstract}
\noindent 监测机构需要知道多收集病例能使再生数估计精确多少。本文在负二项传播与二项报告的框架下区分两个估计目标。对瞬时再生数，报告数除以预计报告规模 $m$ 的估计量在任意规模上达到 Cramér--Rao 下界。对跨期平均再生数，条件无偏估计的误差方差为环境方差 $v_R$、报告异质性方差 $v_H$ 与 $A_H/m$ 之和，前两项不随规模下降。各期环境独立（$\phi=0$）时，这一误差也对应下一期传播水平的一步前瞻预测；均值由此前 $T_0$ 期估计时，误差再乘以 $1+1/T_0$。在此基础上，本文给出提高报告率与扩大覆盖的排序反转阈值、报告噪声在预测误差截距上随滞后权重变化的 $F$ 签名、限制 $v_R$ 本身估计精度的有效期数 $T_{\mathrm{eff}}$，以及环境自相关时由预算阈值决定的期数与规模设计。美国国家医疗安全网络 16,218 个州周的流感住院数据中，流行初期预测误差随规模按 $a+b/D$ 下降，$b/a$ 为 22.6 例/周（95\% 百分位区间 13.5--34.9）。可由各州填报医院比例观测到的报告成分不到截距超出量的 1\%，与填报比例无关的报告离散不受此约束。在流感、COVID-19 与 RSV 的 7 个病原与窗口组合中，4 个部分复现，3 个只呈方向一致，没有组合复现。4 个流行季下季层面方差的相对标准误下界为 0.71，其大小仍无法确定。

\vspace{0.6em}
\noindent\textbf{关键词：}再生数；环境随机性；报告异质性；监测设计；方差分解；流感住院
\end{abstract}
